% MẪU 2 — Đề kiểm tra: PHẦN I trắc nghiệm, PHẦN II đúng sai, PHẦN III trả lời ngắn.
% Lời dặn và số câu mỗi phần tự điền. Bảng đáp án ở cuối chỉ in ở bản GV.
% Đổi lớp: chỉ sửa dòng \usepackage{toan11}.
\documentclass[12pt]{article}
\usepackage{toan11}
\begin{document}
\kythi{Kiểm tra giữa học kì I}
\thoigian{90 phút}
\made{101}
% \so{...} \truong{...}  : chỉ ghi khi khác mặc định trong tritoan.sty
\taodaude

\phantn
\begin{tn}
Tập xác định của hàm số $y=\tan x$ là
\paa{$\mathbb{R}$.}{\dung{$\mathbb{R}\setminus\left\{\frac{\pi}{2}+k\pi \mid k\in\mathbb{Z}\right\}$.}}{$\mathbb{R}\setminus\{k\pi \mid k\in\mathbb{Z}\}$.}{$[-1;1]$.}
\loigiai{$\tan x$ xác định khi $\cos x\ne0$.}
\end{tn}
\begin{tn}
Giá trị của $\sin\dfrac{\pi}{6}$ bằng
\pa{\dung{$\dfrac12$}}{$\dfrac{\sqrt3}{2}$}{$1$}{$0$}
\end{tn}

\phands
\begin{ds}
Cho cấp số cộng $(u_n)$ có $u_1=2$, công sai $d=3$.
\ydung{$u_2=5$.}
\ysai{$u_5=15$.}
\ydung{$S_4=26$.}
\ysai{$(u_n)$ là dãy số giảm.}
\loigiai{$u_5=u_1+4d=14$; $S_4=\frac{4(2u_1+3d)}{2}=26$.}
\end{ds}

\phantln
\begin{tln}
Phương trình $\sin x=0$ có bao nhiêu nghiệm thuộc đoạn $[0;2\pi]$?
\traloi{3}
\end{tln}
\begin{tln}
Cho cấp số nhân có $u_1=3$, $q=2$. Tính $u_4$.
\traloi{24}
\end{tln}

\hetde
\newpage
\bangdapan
\end{document}
